\documentclass[11pt,a4paper]{article}
\usepackage[utf8]{inputenc}
\usepackage[T1]{fontenc}
\usepackage[polish]{babel}
\usepackage{lmodern}
\usepackage[margin=2.3cm]{geometry}
\usepackage{booktabs}
\usepackage{array}
\usepackage{enumitem}
\usepackage{titlesec}
\usepackage{xcolor}
\usepackage{fancyhdr}
\usepackage{hyperref}

\definecolor{navy}{RGB}{27,42,86}
\definecolor{lgray}{RGB}{240,240,244}

\titleformat{\section}{\large\bfseries\color{navy}}{\thesection.}{0.5em}{}
\titleformat{\subsection}{\bfseries\color{navy}}{\thesubsection.}{0.5em}{}

\pagestyle{fancy}
\fancyhf{}
\renewcommand{\headrulewidth}{0.6pt}
\fancyhead[L]{\small Neurofizjologia i pomiar EEG}
\fancyhead[R]{\small Instrukcja laboratoryjna nr 2}
\fancyfoot[C]{\small\thepage}

\newcommand{\infobox}[1]{%
  \par\noindent\colorbox{lgray}{\parbox{\dimexpr\linewidth-2\fboxsep}{\vspace{2pt}#1\vspace{2pt}}}\par\vspace{4pt}}

\title{\vspace{-1.5cm}\textbf{\color{navy}Instrukcja laboratoryjna nr 2}\\[4pt]
\Large Sprzęt do badań EEG i procedura przygotowania rejestracji}
\author{Neurofizjologia i pomiar EEG\\ Psychologia, studia II stopnia --- specjalność neuropsychologia kliniczna}
\date{Czas trwania zajęć: 120 minut}

\begin{document}
\maketitle
\thispagestyle{fancy}

\section{Cele zajęć}
\begin{itemize}[nosep]
  \item Zapoznanie z rodzajami elektrod i czepków EEG oraz zasadą działania wzmacniacza sygnału.
  \item Nauka prawidłowego przygotowania stanowiska badawczego zgodnie z systemem 10--20.
  \item Ćwiczenie procedury zakładania czepka i osiągania niskiej impedancji elektrod.
  \item Kształtowanie właściwego kontaktu z osobą badaną podczas przygotowania do badania.
\end{itemize}

\section{Wprowadzenie teoretyczne}

\subsection{Rodzaje elektrod}
\begin{itemize}[nosep]
  \item \textbf{Elektrody mokre (żelowe/pastowe)} -- Ag/AgCl, wymagają żelu przewodzącego; wysoka jakość sygnału, niska impedancja (docelowo $<5$--$10\,\text{k}\Omega$), czasochłonne przygotowanie.
  \item \textbf{Elektrody suche} -- brak konieczności aplikacji żelu, szybsze zakładanie, wyższa i mniej stabilna impedancja, częściej stosowane w rozwiązaniach mobilnych/konsumenckich.
  \item \textbf{Elektrody aktywne vs. pasywne} -- elektrody aktywne zawierają wstępny wzmacniacz przy samej elektrodzie, co zwiększa odporność na zakłócenia przy dłuższych kablach.
\end{itemize}

\subsection{System rozmieszczenia elektrod 10--20}
Międzynarodowy system 10--20 (Jasper, 1958) definiuje pozycje elektrod na podstawie procentowych odległości (10\% i 20\%) między punktami anatomicznymi: nasion, inion oraz punktami przedusznymi. Litery oznaczają okolicę (Fp -- czołowo-biegunowa, F -- czołowa, C -- centralna, P -- ciemieniowa, O -- potyliczna, T -- skroniowa), a cyfry -- stronę (parzyste: prawa, nieparzyste: lewa, z: linia środkowa). Zagęszczone warianty (10--10, 10--5) stosuje się przy większej liczbie kanałów.

\subsection{Wzmacniacz sygnału i impedancja}
Wzmacniacz EEG podnosi amplitudę mikrowoltowego sygnału (rzędu 10--100 $\mu$V) do poziomu umożliwiającego digitalizację, jednocześnie odfiltrowując zakłócenia sieciowe (50 Hz) dzięki wzmocnieniu różnicowemu (\emph{common mode rejection}). Wysoka i niejednorodna impedancja między elektrodami zniekształca sygnał i zwiększa podatność na artefakty -- dlatego kontrola impedancji jest obowiązkowym etapem przed rejestracją.

\infobox{\textbf{Zasada bezpieczeństwa i komfortu:} Osoba badana musi być informowana o każdym etapie przygotowania, mieć możliwość przerwania procedury oraz komfortowe warunki (temperatura, oświetlenie, brak pośpiechu).}

\section{Sprzęt i materiały}
Czepek EEG (rozmiary S/M/L), zestaw elektrod Ag/AgCl, żel przewodzący, strzykawki do aplikacji żelu, taśma miernicza, wzmacniacz EEG, oprogramowanie do kontroli impedancji, ręczniki papierowe, spirytus do oczyszczenia skóry.

\section{Przebieg zajęć (harmonogram, 120 min)}
\begin{center}
\begin{tabular}{@{}p{2.3cm}p{11.7cm}@{}}
\toprule
\textbf{Czas} & \textbf{Aktywność}\\
\midrule
0--10 min & Powtórzenie materiału z zajęć 1, wprowadzenie do systemu 10--20.\\
10--30 min & Prezentacja rodzajów elektrod, czepków i wzmacniacza -- pokaz sprzętu.\\
30--45 min & Demonstracja prowadzącego: pomiar głowy, dobór rozmiaru czepka, zakładanie czepka.\\
45--55 min & Przerwa.\\
55--90 min & Ćwiczenia praktyczne w parach: zakładanie czepka i aplikacja żelu (zadanie 1).\\
90--110 min & Kontrola impedancji i korekta przygotowania elektrod (zadanie 2).\\
110--120 min & Zdjęcie czepka, omówienie trudności, podsumowanie.\\
\bottomrule
\end{tabular}
\end{center}

\section{Zadania do wykonania}
\begin{enumerate}[label=\textbf{Zadanie \arabic*.}, leftmargin=2.7cm]
  \item W parach: zmierzcie obwód głowy oraz odległość nasion--inion partnera/partnerki, dobierzcie odpowiedni rozmiar czepka i prawidłowo go załóżcie, lokalizując punkt Cz.
  \item Zaaplikujcie żel przewodzący do min. 8 wskazanych przez prowadzącego elektrod, a następnie zmierzcie impedancję każdej z nich, dążąc do wartości poniżej progu ustalonego przez prowadzącego.
  \item Sporządźcie krótki protokół przygotowania badania (checklist) zawierający wszystkie kroki procedury wraz z potencjalnymi źródłami błędu.
\end{enumerate}

\section{Kryteria zaliczenia ćwiczenia}
Poprawność wykonania procedury zakładania czepka i lokalizacji elektrod (40\%), osiągnięta jakość impedancji (30\%), jakość merytoryczna protokołu/checklisty (30\%).

\section{Pytania kontrolne}
\begin{enumerate}[nosep]
  \item Na czym polega procentowy podział głowy w systemie 10--20?
  \item Dlaczego wysoka impedancja pogarsza jakość sygnału EEG?
  \item Jakie są zalety i wady elektrod suchych względem mokrych?
  \item Jakie elementy procedury służą zapewnieniu komfortu osoby badanej?
\end{enumerate}

\section{Literatura}
\begin{enumerate}[nosep]
  \item Szumska, I. (2019). \emph{Neurofizjologia i neuropsychologia -- skrypt dla studentów AEH}. Wydawnictwo AEH.
  \item Biasiucci, A., Franceschiello, B., \& Murray, M. M. (2019). Electroencephalography. \emph{Current Biology, 29}(3), R80--R85.
  \item Kalat, J. W. (2019). \emph{Biological Psychology}. Cengage Learning.
\end{enumerate}

\end{document}
